% =====================================================================
%  CHAPTER 18: 学校突发公共卫生事件应急管理（对应大纲第十一章）
%  参考PPT：第十八章 学校突发公共卫生事件应急管理（李生慧，上海交通大学）
% =====================================================================
\chapter{学校突发公共卫生事件应急管理}
\setcounter{chapter}{18}
\chapterintro{
掌握学校突发公共卫生事件的定义、类型（7种）和分级（I~IV级，对应红橙黄蓝）；掌握学校突发公共卫生事件应对的目标、原则（6项）和措施（政府、教育行政部门、学校三个层面）；掌握学校传染病事件的定义、流行病学特征（地区、时间、病种）和应急管理（疫情报告24h内、调查、控制、信息发布、应急反应终止、评估总结）；掌握学校食源性疾病、食物中毒和学校食物中毒事故的概念；掌握学校食物中毒的类型（细菌性、真菌性、有毒动植物、化学性）、流行病学特征（地区、时间、人群、病因）和应急管理（报告2h内、封存销毁、组织救治、配合调查、恢复教学秩序、信息发布）；掌握群体心因性反应事件的定义、触发因素、个体易患因素、早期识别指标体系和应急管理（应急预案、现场应急处置、心理治疗）。
}

% ----- 12.1 概述 -----
\section{学校突发公共卫生事件概述}

\begin{keybox}
\textbf{学校突发公共卫生事件（Public Health Emergency Events in Schools）：}在学校内突然发生，造成或可能造成师生员工健康严重损害的重大传染病疫情、群体性不明原因疾病、重大食物或职业中毒及其他严重影响师生员工身心健康的公共卫生事件。

\textbf{一、类型（7种）}

\begin{enumerate}[leftmargin=*]
    \item \imp{重大传染病疫情}
    \item \imp{预防接种和预防服药群体性不良反应}
    \item \imp{群体性不明原因疾病}
    \item \imp{食物中毒}
    \item \imp{其他中毒}
    \item \imp{环境因素事件}
    \item \imp{意外辐射照射事件}
\end{enumerate}

\textbf{二、分级（高→低，4级）}

\begin{enumerate}[leftmargin=*]
    \item \imp{特别重大（I级）——红色}
    \item \imp{重大（II级）——橙色}
    \item \imp{较大（III级）——黄色}
    \item \imp{一般（IV级）——蓝色}
\end{enumerate}

\textbf{三、应对}

\textbf{（一）目标：}有效预防、及时控制和消除突发公共卫生事件及其危害，最大程度减少突发公共卫生事件造成的危害，保障师生员工身心健康与生命安全。

\textbf{（二）原则（6项）：}
\begin{enumerate}[leftmargin=*]
    \item \imp{以人为本，减少危害}
    \item \imp{居安思危，预防为主}
    \item \imp{统一领导，分级负责}
    \item \imp{依法规范，加强管理}
    \item \imp{快速反应，协同应对}
    \item \imp{依靠科技，提高素质}
\end{enumerate}

\textbf{（三）措施（三个层面）}

\textbf{1. 政府层面：}
\begin{itemize}[leftmargin=*]
    \item 建立应急组织和指挥机构
    \item 建立专家咨询委员会
    \item 建立应急处理专业技术机构
\end{itemize}

\textbf{2. 教育行政部门层面：}
\begin{itemize}[leftmargin=*]
    \item 建立组织机构
    \item 制订应急预案
    \item 应急反应
    \item 信息发布
\end{itemize}

\textbf{3. 学校层面：}
\begin{itemize}[leftmargin=*]
    \item 加强领导，实行单位领导负责制
    \item 健全学校各项健康管理制度，加强学生健康管理
    \item 认真落实突发公共卫生事件报告管理制度，确保报告信息畅通
    \item 配合卫生部门采取有效措施，及时控制事件发展蔓延
    \item 加强健康教育，提高学生应对能力
\end{itemize}
\end{keybox}

% ----- 12.2 传染病事件 -----
\section{学校传染病事件应对}

\begin{keybox}
\textbf{学校传染病事件（Infectious Disease Events in School）：}学校某种传染病在短时间内发生、波及范围广泛，出现大量病人或死亡病例，发病率远超常年发病率水平的情况。

\textbf{一、流行病学特征}

学校突发公共卫生事件最主要类型。

\begin{enumerate}[leftmargin=*]
    \item \imp{地区分布}：社会经济相对落后的西南省份、乡村中小学和县城小学多发。
    \item \imp{时间分布}：双峰型，3~6月、10~12月。
    \item \imp{病种分布}：呼吸道（流感、水痘、风疹、麻疹、流行性腮腺炎）和消化道传染病为主。
\end{enumerate}
\textbf{二、应急管理}

\textbf{1. 疫情报告（24h内）：}发生法定传染病疫情或突发公共卫生事件时，学校疫情报告人应在24小时内向属地疾病预防控制机构和教育行政部门报告。报告依据：在同一宿舍或者同一班级，1天内有3例或者连续3天内有多个学生（5例以上）患病，并有相似症状（如发热、皮疹、腹泻、呕吐、黄疸等）或者共同用餐、饮水史；个别学生出现不明原因的高热、呼吸急促或剧烈呕吐、腹泻等症状；学校发生群体性不明原因疾病或者其他突发公共卫生事件。

\textbf{2. 疫情调查：}学校全面配合卫生部门开展疫情分析、病例诊治以及流行病学调查和疫情处理工作。

\textbf{3. 控制疫情：}
\begin{itemize}[leftmargin=*]
    \item 及时隔离，安排就诊：对确诊患有法定传染病的学生、疑似病人或传染病密切接触者，配合卫生部门依法进行隔离或者医学观察，并及时安排就诊，做好检疫期相关记录。
    \item 疫情防控合作：配合属地疾病预防控制机构开展消毒、疫情调查和宣传教育等工作。
    \item 应急疫苗接种与预防用药：根据需要组织开展应急疫苗接种、预防性服药。
    \item 复课：学生病愈且隔离期满时，应到学校医务室或卫生室查验后方可进班复课。
    \item 停课并暂停集会：在传染病暴发、流行时，学校应停止举办大型师生集会和会议，采取临时停课或暂时关闭措施。
    \item 做好个人防护：教职员工在照顾患病学生、接触可能受到污染的物品或排泄物时，应根据实际情况采取必要的个人防护措施。
\end{itemize}

\textbf{4. 信息发布：}教育部门和学校要按照有关规定作好信息发布工作，信息发布要及时主动、实事求是，准确把握、注意技巧，正确引导舆论，维护社会稳定。

\textbf{5. 健康教育：}学校应根据事件性质，有针对性地开展相关传染病预防知识的教育活动，提高师生健康意识和自我防护能力，开展心理危机干预工作以消除学生的情绪波动和心理障碍。

\textbf{6. 应急反应终止：}学校传染病事件应急反应终止的条件是末例传染病病例发生后经过最长潜伏期无新的病例出现。由卫生行政部门对学校传染病事件作出应急反应的终止决定后，学校方可解除应急措施。

\textbf{7. 评估总结：}学校在传染病暴发、流行事件得到控制后，对事件进行评估总结，评估内容主要包括事件概况、现场调查处理概况、病人救治情况、所采取措施的效果评价、应急处理过程中存在的问题和取得的经验及改进建议。评估报告须上报教育行政部门。
\end{keybox}

% ----- 12.3 食物中毒 -----
\section{学校食物中毒事件应对}

\begin{keybox}
\textbf{一、概念}

\textbf{1. 食源性疾病（Food Origin Disease）：}食品中致病因素进入人体引起的感染性、中毒性疾病。3个基本要素：
\begin{itemize}[leftmargin=*]
    \item \imp{致病因子}：病毒、细菌、寄生虫、毒素、重金属、有毒有害化学物质
    \item \imp{传播媒介}：食物
    \item \imp{临床表现/症状体征}：轻微胃肠炎→致命神经毒作用、肝肾综合征等各不相同
\end{itemize}

\textbf{2. 食物中毒（Food Poisoning）：}食用被有毒有害物质污染的食品或食用含有毒有害物质的食品后出现的急性、亚急性疾病。特点：
\begin{itemize}[leftmargin=*]
    \item 发病与特定食物有关
    \item 潜伏期短，发病曲线单峰型
    \item 临床表现相似：最常见为消化道症状，如腹痛、腹泻、恶心、呕吐
    \item 一般没有人与人间的直接传染
\end{itemize}

\textbf{3. 学校食物中毒事故（Food Poisoning Incidents in School）：}由学校主办或管理的校内供餐单位及学校负责组织提供的集体用餐导致的学校师生食物中毒事故，分为重大、较大和一般学校食物中毒事故。
\textbf{二、流行病学特征}

学校突发公共卫生事件主要类型，仅次于传染病。

\begin{enumerate}[leftmargin=*]
    \item \imp{地区}：南方$>$北方，华东、华南最高发。
    \item \imp{时间}：秋季最高，夏季次之，冬季最少。
    \item \imp{人群}：(1)食物中毒起数、中毒人数：中学$>$小学；(2)死亡病例集中在小学，尤其是农村小学。
    \item \imp{病因}：食物污染和变质是最常见原因，微生物是最主要致病因子，化学物是引起死亡的主要因素。
\end{enumerate}
\textbf{三、应急管理}

\begin{enumerate}[leftmargin=*]
    \item \imp{报告（2h内）}：2小时内报告当地教育行政部门和卫生行政部门；报告的主要内容包括发生食物中毒暴发事件单位、地点、时间、中毒人数、主要临床症状等。
    \item \imp{封存、销毁}：对疑似导致食品安全事故的食品及其原料立即封存，并由卫生机构进行检验；对确认属于被污染的食品及其原料，立即停止销售，予以召回和销毁。封存被污染的食品用工具及用具，并进行彻底清洗消毒。
    \item \imp{迅速组织救治}：发生学校食物中毒事故后，教育行政部门和学校要把治病救人放在首位，迅速组织救治，尤其是危重患者，要不惜一切代价，全力抢救。
    \item \imp{配合调查和善后}：学校应配合卫生行政部门进行调查，按卫生行政部门的要求如实提供有关材料和样品。
    \item \imp{尽早恢复正常教学秩序}：学校食物中毒事故应急处置完成后，学校应尽快恢复正常的教学秩序。
    \item \imp{正常发布相关信息}：信息发布要及时主动、准确把握，实事求是，正确引导舆论，注重社会效果。
\end{enumerate}
\end{keybox}

% ----- 12.4 群体心因性反应 -----
\section{学校群体心因性反应事件应对}

\begin{keybox}
\textbf{一、群体心因性反应（Mass Psychogenic Reaction）}

一种群体精神性反应，是在一定社会文化背景条件下，在两人或两人以上群体中发生、有躯体性疾病症候群、但没有可检测出的器质性变化的病症。

\textbf{1. 触发因素：}公共卫生事件为主。其中尤以集体预防接种和服药、疑似食物中毒最为常见，而迷信谣言、心理压力和情绪紧张等因素也常有报道。

\textbf{2. 个体易患因素：}
\begin{itemize}[leftmargin=*]
    \item 农村多发
    \item 小学和初中多发
    \item 女生$>$男生
    \item 认知能力较差，易接受心理暗示
\end{itemize}

\textbf{3. 临床分型：}
\begin{enumerate}[leftmargin=*]
    \item \imp{焦虑型癔症（Anxiety Hysteria）}：常见症状包括腹痛、头痛、头晕、晕厥、恶心和换气过度。20世纪以来最为常见，常由突然接触或目击引发焦虑的情景（如异常气味、疑似食物中毒）触发，以焦虑相关躯体化表现为主。
    \item \imp{运动型癔症（Motor Hysteria）}：常见症状包括舞蹈样动作、抽搐、颤抖、行走困难、无法控制的大笑和哭泣、沟通困难和恍惚状态。20世纪前为主流类型，主要由宗教迷信导致的长期压抑所诱发，以运动功能障碍和异常行为表现为特征。
\end{enumerate}
\imp{鉴别要点：}焦虑型以感觉和自主神经功能障碍为主，运动型以运动功能障碍为主；两型均无器质性病变证据。
\end{keybox}

\subsection{早期识别}

\begin{keybox}
\textbf{群体心因性反应事件的早期识别指标体系：}

\begin{enumerate}[leftmargin=*]
    \item \imp{症状特点}：突然出现的集体性躯体症状（头痛、头晕、腹痛、恶心、呕吐、过度换气、晕厥等），症状多样但均缺乏器质性病变证据。

    \item \imp{时间聚集性}：短时间内（数分钟至数小时）多人相继出现类似症状，发病曲线呈爆发性上升。

    \item \imp{无器质性病变}：体格检查、实验室检查和影像学检查均无阳性发现，症状与客观检查结果严重不符。

    \item \imp{人群特征}：女性多于男性，小学生和初中生高发，认知能力较差、易接受心理暗示的个体更易受影响。

    \item \imp{传播模式}：症状通过视觉、听觉等感官暗示途径在人群中传播，首发病例（index case）往往是群体关注的焦点。

    \item \imp{情绪关联性}：症状的发生和严重程度与情绪紧张、恐惧、焦虑等心理因素密切相关，在消除紧张情绪环境后症状可迅速缓解。

    \item \imp{社会心理背景}：常发生在集体预防接种/服药、疑似食物中毒、重要考试等情绪高度紧张的事件背景下。
\end{enumerate}

\textbf{识别要点：}早期原因不明时，需综合上述指标进行判断，特别重视"无器质性病变"这一核心鉴别特征，同时排除器质性疾病（如真正的食物中毒、传染病暴发等）的可能。
\end{keybox}

\textbf{三、应急管理}

\subsection{应急预案}

\begin{keybox}
\textbf{群体心因性反应事件的应急预案应包括以下内容：}

\begin{enumerate}[leftmargin=*]
    \item \imp{组织指挥体系}：成立由学校领导、校医、心理教师、班主任等组成的应急指挥小组，明确各成员职责分工，建立与卫生健康部门、教育行政部门的联络机制。

    \item \imp{事件分级响应}：根据受影响人数、症状严重程度、波及范围等将事件分为不同响应等级，制定对应级别的处置措施和资源调配方案。

    \item \imp{现场处置流程}：
    \begin{itemize}
        \item 隔离患者，转移出现场进行适当隔离，避免互相影响及效仿
        \item 消除紧张性情绪环境，引导其他学生撤离，消除不良言语/动作暗示
        \item 由医生进行详细检查，向患者解释检查结果，帮助建立信心
        \item 症状缓解后帮助分析主客观原因，进行心理疏导
    \end{itemize}

    \item \imp{信息报告与沟通}：明确向上级部门报告的程序、时限和内容要求，制定向家长和社会公众的信息通报策略，统一信息发布口径，防止谣言传播和恐慌扩散。

    \item \imp{心理干预方案}：制定分级心理干预措施（一般性安抚$\rightarrow$团体心理辅导$\rightarrow$个体心理咨询$\rightarrow$专业转介），明确心理教师和校外心理援助资源的调用程序。

    \item \imp{后续处置与评估}：事件平息后的跟踪随访计划、心理康复评估、健康教育强化措施，以及事件处置的总结评估和改进建议。
\end{enumerate}
\end{keybox}

\begin{keybox}
\textbf{（二）现场应急处置：}
\begin{enumerate}[leftmargin=*]
    \item \imp{隔离患者}：立即将首发、继发病例转移出现场，适当隔离，以避免患者间互相影响及效仿，减少症状的顽固性和丰富性。
    \item \imp{消除紧张性情绪环境}：(1)在隔离患者的同时，引导其他学生也从该环境撤离；(2)消除来自周围环境的不良的言语暗示、动作暗示；(3)医生认真做详细检查，解释病情，使学校领导、教师、家属和围观者对治疗建立信心；(4)待症状缓解后，帮助患者分析发病的主客观原因，指导他们和家属解除相关不良精神因素。
    \item \imp{对症治疗}：对患者表现出的躯体症状采用止吐、止泻、止痛等对症治疗；对少数精神反应特别大的个体可适量使用镇静药物。
\end{enumerate}

\textbf{（三）心理治疗：}
\begin{enumerate}[leftmargin=*]
    \item \imp{移情解释法}：采用符合患者心理要求的内容及形式，鼓励其参加感兴趣的活动以转移注意力。
    \item \imp{暗示疗法}：有语言暗示、药物暗示、催眠疗法等。
    \item \imp{着重治疗关键患者（首发患者）}：通过重点治疗使其痊愈，再通过其现身说法，对其他患者产生正向影响，对控制事件发展起事半功倍的作用。
    \item \imp{争取家长配合}：学校群体心因性疾病的触发因素并非直接影响儿童，而是经过家长等成人的过滤、整合后传递给儿童成为暗示信号。在预防、应急处理和后续的心理干预阶段均应重视对家长的教育引导解释。
\end{enumerate}
\end{keybox}

\newpage
